% Capitulo 2: Antecedentes (Proyecto de Tesis -- Fase 8, docs/rio_search_plan.md S3.11)
\chapter{Antecedentes}
\label{chap:antecedentes}

Las redes \emph{Long Short-Term Memory} \citep{hochreiter-1997-lstm} resuelven el desvanecimiento
del gradiente que limita a las redes recurrentes cl\'asicas en secuencias largas, mediante una
celda de memoria con compuertas de entrada, olvido y salida -- la raz\'on por la que se propone una
variante bidireccional como primer modelo no trivial de este proyecto.
\citep{kratzert-2018-lstm-rainfall-runoff} muestran que las LSTM igualan o superan a modelos
hidrol\'ogicos conceptuales calibrados por cuenca cuando se entrenan sobre variables meteorol\'ogicas
de entrada sin par\'ametros f\'isicos ajustados a mano, lo que motiva un enfoque orientado a datos
en vez de un modelo hidrol\'ogico conceptual expl\'icito.

Para la evaluaci\'on, \citep{gupta-2009-kge} descomponen el error cuadr\'atico medio normalizado en
correlaci\'on, sesgo en la media y sesgo en la variabilidad, y proponen el \emph{Kling-Gupta
Efficiency} (KGE) como m\'etrica que pondera esos tres componentes por separado -- la m\'etrica que
este proyecto adopta para seleccionar el modelo campe\'on sobre el conjunto de validaci\'on.

El formato de este documento se extrajo de un trabajo final previo del autor, de la misma carrera,
disponible como modelo en \texttt{rio\_search/research/templates/}: comparte el enfoque general de
predicci\'on sobre series temporales hidrol\'ogicas con aprendizaje autom\'atico, pero no su objeto
de estudio.
